\chapter{Construction and shared ownership}
\label{chap:creators}

The public constructors create values. The allocation helpers in
\api{nurbspath/creators.hpp} create the corresponding values under
\api{std::shared_ptr} ownership. They retain the constructors' geometric
validation and exceptions. Both styles are useful: local calculations often
use values, while a collection with shared lifetime may use factory results.
The umbrella header includes these helpers.

\begin{lstlisting}[language=C++]
#include <nurbspath/config.hpp>
#include <nurbspath/nurbspath.hpp>

using namespace nurbspath;
point3<double> value{1.0, 2.0, 3.0};
auto shared = make_point3<double>(1.0, 2.0, 3.0);
const double separation = distance(value, *shared);
\end{lstlisting}

The suffix identifies the Cartesian world. A 2D factory never silently
allocates a 3D entity. Calls below use \api{double}; replacing it with
another standard floating-point type changes the scalar type throughout.
Allocation can throw \api{std::bad_alloc} in addition to the validation
exceptions noted for each entity.

\section{Point and vector factory catalogue}

\begin{longtable}{@{}p{.43\linewidth}p{.52\linewidth}@{}}
\caption{All point and vector allocation families. Each result is a shared
pointer to the type named after \api{make_}.}\\
\toprule
Call arguments after the template & Created value \\
\midrule
\endfirsthead
\toprule Call arguments after the template & Created value \\
\midrule
\endhead
\api{make_point2<REAL>()} & The 2D origin. \\
\api{make_point2<REAL>(x, y)} & A 2D position with the supplied coordinates. \\
\api{make_point3<REAL>()} & The 3D origin. \\
\api{make_point3<REAL>(x, y, z)} & A 3D position with the supplied coordinates. \\
\api{make_vector2<REAL>()} & The 2D zero vector. \\
\api{make_vector2<REAL>(x, y)} & The supplied 2D components. \\
\api{make_vector2<REAL>(point)} & Coordinates copied from a \api{point2}. \\
\api{make_vector2<REAL>(circle)} & Center coordinates copied from a \api{circle2}. \\
\api{make_vector2<REAL>(ray)} & Origin coordinates copied from a \api{ray2}. \\
\api{make_vector3<REAL>()} & The 3D zero vector. \\
\api{make_vector3<REAL>(x, y, z)} & The supplied 3D components. \\
\api{make_vector3<REAL>(point)} & Coordinates copied from a \api{point3}. \\
\api{make_vector3<REAL>(sphere)} & Center coordinates copied from a \api{sphere3}. \\
\api{make_vector3<REAL>(plane)} & Parameter-origin coordinates copied from a \api{plane3}. \\
\api{make_vector3<REAL>(ray)} & Origin coordinates copied from a \api{ray3}. \\
\bottomrule
\end{longtable}

The entity-to-vector factories perform the same explicit component copy as
the constructors in Chapter~\ref{chap:entities}. To allocate a copy of a
ray direction, obtain \api{ray.direction()} and use its components, or use
\api{std::make_shared<vector3<double>>(ray.direction())}.

\section{Ray, circle, and sphere creation}

The origin-and-direction helpers are \api{make_ray2} and \api{make_ray3}.
The two-point helpers are \api{make_ray2_from_points} and
\api{make_ray3_from_points}.
The first takes \api{(origin, direction, tolerance)}; the second takes
\api{(origin, through_point, tolerance)}. Tolerance is optional and defaults
to \api{vector2<REAL>::default_tolerance()} or its 3D counterpart.
The point-based helper uses \api{through_point - origin} directly, with
no normalization. The through point is therefore at native parameter $s=1$.

\begin{lstlisting}[language=C++]
const point2<double> a2{1.0, 1.0};
const point2<double> b2{3.0, 2.0};
auto ray_a = make_ray2<double>(a2, b2 - a2);
auto ray_b = make_ray2_from_points<double>(a2, b2);
auto circle = make_circle2<double>(point2<double>{0.0, 0.0}, 2.0);

const point3<double> a3{1.0, 1.0, 0.0};
const point3<double> b3{3.0, 2.0, 1.0};
auto ray_c = make_ray3<double>(a3, b3 - a3);
auto ray_d = make_ray3_from_points<double>(a3, b3);
auto sphere = make_sphere3<double>(point3<double>{0.0, 0.0, 0.0}, 2.0);
\end{lstlisting}

A direction or point separation within the ray tolerance throws
\api{std::invalid_argument}. \api{make_circle2(center, radius)} and
\api{make_sphere3(center, radius)} require finite, positive radii and throw
\api{std::invalid_argument} otherwise. These are the complete circle and
sphere factory families; they do not infer a circle or sphere from several
boundary points.

\section{Five ways to allocate a plane}

The three \api{make_plane3} overloads mirror the three constructors:

\begin{enumerate}[leftmargin=*]
\item \api{(normal, signed_distance, tolerance)} uses Hessian normal form.
\item \api{(origin, normal, tolerance)} chooses a stable in-plane frame.
\item \api{(origin, normal, u_hint, tolerance)} sets the preferred positive-u
direction by in-plane projection.
\end{enumerate}

Two named helpers derive an oriented frame from construction geometry:
\begin{itemize}[leftmargin=*]
\item \api{make_plane3_from_u_direction} takes \api{origin},
\api{u_direction}, \api{plane_point}, and optional \api{tolerance}.
\item \api{make_plane3_from_points} takes \api{origin}, \api{u_point},
\api{plane_point}, and optional \api{tolerance}; it uses
\api{u_point - origin} as its positive-u direction.
\end{itemize}
The third point determines the positive-v side by rejection from that axis:
\[
 \widehat u=\frac{U}{\lVert U\rVert},\quad
 V=(P-O)-((P-O)\cdot\widehat u)\widehat u,\quad
 \widehat v=\frac{V}{\lVert V\rVert},\quad
 \widehat n=\widehat u\times\widehat v.
\]
All five helpers accept the optional vector tolerance. A short normal,
u-direction, or rejected v-component causes a \api{std::domain_error}
exception. A non-finite Hessian distance causes \api{std::invalid_argument}.

\begin{lstlisting}[language=C++]
const point3<double> origin{1.0, 2.0, 3.0};
const point3<double> u_point{3.0, 2.0, 3.0};
const point3<double> plane_point{2.0, 4.0, 3.0};
auto hessian = make_plane3<double>(vector3<double>{0.0, 0.0, 2.0}, 3.0);
auto from_normal = make_plane3<double>(origin, vector3<double>::unit_z());
auto with_hint = make_plane3<double>(
    origin, vector3<double>::unit_z(), vector3<double>::unit_x());
auto from_direction = make_plane3_from_u_direction<double>(
    origin, vector3<double>{2.0, 0.0, 0.0}, plane_point);
auto from_points = make_plane3_from_points<double>(origin, u_point, plane_point);
\end{lstlisting}

\begin{figure}[htbp]
\centering
\begin{tikzpicture}[scale=1.1]
\fill[surface] (-.3,-.35)--(4.7,-.35)--(4.7,2.75)--(-.3,2.75)--cycle;
\coordinate (O) at (0,0); \coordinate (U) at (3.8,0);
\coordinate (P) at (2.2,2.3); \coordinate (H) at (2.2,0);
\draw[vector] (O)--(U) node[midway,below]{$U-O$};
\draw[guide] (O)--(P);
\draw[guide] (P)--(H);
\draw[vector,manualteal] (H)--(P) node[midway,right]{$V$};
\draw[vector,manualorange] (O)--(0,1.15) node[left]{$\widehat v$};
\node[point,label=below left:$O$] at (O) {};
\node[point,label=below:$U$] at (U) {};
\node[point,label=above:$P$] at (P) {};
\draw[manualmuted] (2,.0)--(2,.2)--(2.2,.2);
\end{tikzpicture}
\caption{Three-point plane construction viewed in the plane itself. The
third point chooses the positive-v side; its along-u component is discarded
when the frame is built.}
\end{figure}

Moving the third point across the u-axis reverses both $\widehat v$ and
$\widehat n$. Three collinear points do not determine a plane through this
interface. A supplied normal or u-direction need not have unit length.

\section{Ordinary spline constructors and factories}

The following catalogue applies independently to \api{nurbs_spline2} and
\api{nurbs_spline3}. Write \api{S} for either type,
\api{P} for its same-world point type, \api{A} for the corresponding
array-backed type, \api{p} for the degree, \api{n} for the control-point
count, and \api{tol} for the positive definition tolerance. The factories
\api{make_nurbs_spline2} and \api{make_nurbs_spline3} return
\api{std::shared_ptr<S>}.

\begin{table}[htbp]
\centering
\caption{Complete ordinary-spline construction families. Optional brackets
describe permitted omitted arguments, not C++ syntax.}
\begin{tabularx}{\linewidth}{@{}p{.52\linewidth}Y@{}}
\toprule
Argument sequence & Availability and behavior \\
\midrule
\api{start, end, n, p} & Constructor and factory; standard open definition. \\
\api{start, end, n, p, tol} & Constructor and factory; explicit tolerance. \\
\api{start, end, n, p, closed [, tol]} & Constructor and factory; explicit Boolean closure request. \\
\api{points, weights, knots, p [, tol]} & Constructor and factory with \api{std::vector} collections. \\
\api{points, weights, knots, p, closed [, tol]} & Constructor and factory with vector collections and a closure request. \\
\api{points, weights, knots, p [, tol]} & Factory also accepts point/scalar \api{std::valarray} collections. \\
\api{points, weights, knots, p, closed [, tol]} & Factory also accepts valarrays and a closure request. \\
\api{array_spline} & Explicit constructor and factory; copy from corresponding \api{A}. \\
\bottomrule
\end{tabularx}
\end{table}

The omitted spline tolerance is
\api{64 * std::numeric_limits<REAL>::epsilon()}.
The standard constructor places $n$ control points uniformly from start to
end with unit weights and open-clamped knots over
$[0,\operatorname{distance}(\mathrm{start},\mathrm{end})]$. It creates a
straight curve, with a positive domain extent required by validation.
The explicit-tolerance fifth argument is a scalar; the closure overload
requires an actual \api{bool}, avoiding accidental interpretation of a
tolerance as a closure flag.

\begin{lstlisting}[language=C++]
const point2<double> start{0.0, 0.0};
const point2<double> end{4.0, 0.0};
nurbs_spline2<double> line(start, end, 5, 2);
auto shared_line = make_nurbs_spline2<double>(start, end, 5, 2);
auto explicit_tolerance = make_nurbs_spline2<double>(start, end, 5, 2, 1e-10);
auto explicit_open = make_nurbs_spline2<double>(start, end, 5, 2, false, 1e-10);
\end{lstlisting}

Complete definitions supply every control point, positive weight, and knot.
For $n$ control points and degree $p$, the knot count is $n+p+1$, with
$0<p<n$, nondecreasing finite knots, and a nonempty active domain. The
standard and complete-definition families have matching validation contracts;
the full mathematical requirements are in Chapter~\ref{chap:splines}.

\begin{lstlisting}[language=C++]
std::vector<point3<double>> points{
    {0.0, 0.0, 0.0}, {1.0, 2.0, 0.5}, {3.0, 0.0, 1.0}};
std::vector<double> weights{1.0, 1.0, 1.0};
std::vector<double> knots{0.0, 0.0, 0.0, 2.0, 2.0, 2.0};
nurbs_spline3<double> value_curve(points, weights, knots, 2);
auto shared_curve = make_nurbs_spline3<double>(points, weights, knots, 2);

std::valarray<point3<double>> array_points{points.data(), points.size()};
std::valarray<double> array_weights{weights.data(), weights.size()};
std::valarray<double> array_knots{knots.data(), knots.size()};
auto copied_arrays = make_nurbs_spline3<double>(
    array_points, array_weights, array_knots, 2, false);
\end{lstlisting}

The ordinary constructors take vectors. The ordinary factories additionally
accept const, mutable, lvalue, or rvalue valarrays of same-world points plus
scalar valarrays; they copy them in index order into ordinary vector storage.
Brace-list calls remain available through the vector overloads.

\subsection{Closure describes a seam}

A \api{closed = true} request requires coincident evaluated endpoints of
the active domain within the spline tolerance. It does not append a control
point or synthesize a loop from a straight start/end construction. Use a
complete definition or closed interpolation samples to create a useful closed
shape.

\begin{lstlisting}[language=C++]
auto closed_polygon = make_nurbs_spline2<double>(
    std::vector<point2<double>>{
        {0.0, 0.0}, {2.0, 0.0}, {1.0, 1.5}, {0.0, 0.0}},
    std::vector<double>{1.0, 1.0, 1.0, 1.0},
    std::vector<double>{0.0, 0.0, 1.0, 2.0, 3.0, 3.0},
    1, true);
\end{lstlisting}

\begin{figure}[htbp]
\centering
\begin{tikzpicture}[scale=1.1]
\begin{scope}
\draw[curve] (0,0)--(4,0);
\foreach \x in {0,1,2,3,4}{\node[point] at (\x,0) {};}
\node[below] at (0,0) {start}; \node[below] at (4,0) {end};
\node[above] at (2,.35) {standard straight definition};
\end{scope}
\begin{scope}[shift={(6,0)}]
\draw[curve] (0,0)--(2,0)--(1,1.5)--cycle;
\node[point,label=below left:seam] at (0,0) {};
\node[point] at (2,0) {}; \node[point] at (1,1.5) {};
\node[above] at (1,1.9) {complete closed definition};
\end{scope}
\end{tikzpicture}
\caption{The standard factory places control points along a segment. A
closed flag validates a seam; the supplied definition determines the shape.}
\end{figure}

\section{Array-backed spline construction}

\api{nurbs_arr_spline2} and \api{nurbs_arr_spline3} are separate classes,
not derived ordinary splines. Their public \api{std::valarray<REAL> parameters}
contains point-major control coordinates, then all weights, then all knots.
The degree, tolerance, and closure state are stored separately. This layout
is useful when an optimizer must copy or replace the entire scalar state.
The conversion operations retain the Cartesian world.

For dimension $d$, $n$ controls, and degree $p$, the flattened length is
$(d+2)n+p+1$. In 2D it is
\[
 [x_0,y_0,\ldots,x_{n-1},y_{n-1},
   w_0,\ldots,w_{n-1},k_0,\ldots,k_{n+p}],
\]
and in 3D each point contributes X/Y/Z coordinates.

\begin{table}[htbp]
\centering
\caption{Complete array-backed construction families, for both dimensions.
Direct constructors and shared-pointer factories support the same families.}
\begin{tabularx}{\linewidth}{@{}p{.57\linewidth}Y@{}}
\toprule
Argument sequence & Created definition \\
\midrule
\api{start, end, n, p} & Standard open straight curve. \\
\api{start, end, n, p, tol} & Standard curve with explicit tolerance. \\
\api{start, end, n, p, closed [, tol]} & Standard curve with a Boolean closure request. \\
\api{parameters, p [, tol]} & Open definition from flattened scalar valarray. \\
\api{parameters, p, closed [, tol]} & Flattened definition with a Boolean closure request. \\
\api{points, weights, knots, p [, tol]} & Open definition from component vectors. \\
\api{points, weights, knots, p, closed [, tol]} & Component-vector definition with a Boolean closure request. \\
\api{ordinary_spline} & Independent copy from corresponding ordinary spline. \\
\bottomrule
\end{tabularx}
\end{table}

The matching factory names are \api{make_nurbs_arr_spline2} and
\api{make_nurbs_arr_spline3}.
The component collection overloads take \api{std::vector}s; a flattened
overload takes one scalar \api{std::valarray}. The point-valarray overloads
listed for ordinary factories are not separate array-backed factory
overloads. The constructors and factories reject malformed array shapes
before indexing the data.

\begin{lstlisting}[language=C++]
std::valarray<double> parameters{
    0.0, 0.0,  1.0, 2.0,  3.0, 0.0,  // three 2D controls
    1.0, 1.0, 1.0,                    // weights
    0.0, 0.0, 0.0, 2.0, 2.0, 2.0};  // quadratic knots
nurbs_arr_spline2<double> array_curve(parameters, 2);
auto shared_array = make_nurbs_arr_spline2<double>(parameters, 2, false);
\end{lstlisting}

\section{Conversions and definition updates}

Conversions between the ordinary and array-backed representations are
explicit, copying control points, weights, knots, degree, tolerance, closure,
and the native parameter domain. They allocate independent storage:

\begin{lstlisting}[language=C++]
auto ordinary = make_nurbs_spline2<double>(*shared_array);
auto flattened = make_nurbs_arr_spline2<double>(*ordinary);
nurbs_spline2<double> ordinary_value(*flattened);
nurbs_arr_spline2<double> array_value(ordinary_value);
\end{lstlisting}

The same calls are available for the 3D pair. These conversions never add,
remove, or reinterpret a coordinate dimension. Converting an array-backed
definition validates its current public storage; a malformed state throws
rather than yielding a partially constructed spline.

Ordinary spline definition storage is private. Checked setters preserve the
degree, validate a complete proposed state, and refresh cached endpoints only
after success. A failed update leaves the spline unchanged. The array-backed
alternatives expose scalar storage and derive endpoint state when needed;
direct edits must be followed by \api{validate()} or a checked geometry
operation. Chapter~\ref{chap:splines} gives the complete editing and
evaluation APIs for both representations.
